\section{Ceva's theorem and Menelaus' theorem}

Both theorems become one-line determinant criteria in barycentric coordinates.  This is preferable to proofs by areas because the converses and all signs are automatic.

\subsection{Barycentric setup}

Normalize the reference triangle projectively as
\[
A=(1,0,0),\qquad B=(0,1,0),\qquad C=(0,0,1).
\]
Let
\[
D=(0,d_B,d_C)\in BC,\quad
E=(e_A,0,e_C)\in CA,\quad
F=(f_A,f_B,0)\in AB,
\]
where every displayed coordinate is nonzero.  Then the directed side ratios are
\[
\frac{BD}{DC}=\frac{d_C}{d_B},\qquad
\frac{CE}{EA}=\frac{e_A}{e_C},\qquad
\frac{AF}{FB}=\frac{f_B}{f_A}.
\]
These equalities are definitions of barycentric side coordinates and can be kept cross-multiplied.

\subsection{Ceva}

\begin{theorem}[Constructive Ceva]
The cevians $AD,BE,CF$ are concurrent if and only if
\[
d_Ce_Af_B=d_Be_Cf_A.
\]
Equivalently,
\[
\frac{BD}{DC}\frac{CE}{EA}\frac{AF}{FB}=1.
\]
\end{theorem}

\begin{proof}
The three cevian lines have covectors
\[
AD=(0,-d_C,d_B),\qquad
BE=(e_C,0,-e_A),\qquad
CF=(-f_B,f_A,0).
\]
Their determinant is
\[
\det\begin{pmatrix}
0&-d_C&d_B\\
e_C&0&-e_A\\
-f_B&f_A&0
\end{pmatrix}
=d_Be_Cf_A-d_Ce_Af_B.
\]
Thus concurrency implies the displayed product identity.

For the converse, assume the identity and define the explicit point
\[
P=(e_Ad_C,\;d_Be_C,\;d_Ce_C).
\]
It is nonzero.  Direct substitution gives
\begin{align*}
AD\cdot P&=-d_Cd_Be_C+d_Bd_Ce_C=0,\\
BE\cdot P&=e_Ce_Ad_C-e_Ad_Ce_C=0,\\
CF\cdot P&=-f_Be_Ad_C+f_Ad_Be_C=0,
\end{align*}
where the final equality is exactly the hypothesis.  Hence all three cevians pass through $P$.
\end{proof}

This proof supplies the concurrency witness rather than deriving its existence indirectly.  It is therefore intuitionistically stronger than a bare determinant-zero argument.

\subsection{Menelaus}

\begin{theorem}[Constructive Menelaus]
The points $D,E,F$ are collinear if and only if
\[
d_Be_Cf_A+d_Ce_Af_B=0.
\]
Equivalently, in directed ratios,
\[
\frac{BD}{DC}\frac{CE}{EA}\frac{AF}{FB}=-1.
\]
When exactly one of $D,E,F$ lies on an extended side outside the triangle, taking ordinary positive lengths removes the minus sign and gives the customary MathWorld equality of the two triple products.
\end{theorem}

\begin{proof}
The determinant of the three point vectors is
\[
\det\begin{pmatrix}
0&e_A&f_A\\
d_B&0&f_B\\
d_C&e_C&0
\end{pmatrix}
=d_Be_Cf_A+d_Ce_Af_B.
\]
Hence collinearity is equivalent to the stated polynomial equation.

For a completely explicit converse, the line through $D$ and $E$ is
\[
D\join E=(d_Be_C,\;d_Ce_A,\;-d_Be_A).
\]
Its value at $F$ is
\[
(d_Be_C)f_A+(d_Ce_A)f_B,
\]
which vanishes precisely under the Menelaus identity.  Thus $F$ lies on $DE$.
\end{proof}

\subsection{Relation between the two signs}

Ceva and Menelaus differ only in whether the relevant $3\times3$ matrix contains line covectors or point vectors.  In the first determinant the two contributing monomials have opposite signs; in the second they have the same sign.  This is the coordinate source of the familiar $+1$ versus $-1$ distinction in directed-ratio formulations.
